% Part: turing-machines
% Chapter: machines-computations
% Section: combining-machines

\documentclass[../../../include/open-logic-section]{subfiles}

\begin{document}

\olfileid{tur}{mac}{cmb}
\olsection{ترکیب ماشین‌های تورینگ}

\begin{explain}
نمونه‌های ماشین تورینگی که تا اینجا دیده‌ایم نسبتاً ساده بوده‌اند.
اما در واقع هر مسئله‌ای که با هر زبان برنامه‌نویسی امروزی حل‌شدنی است
با ماشین‌های تورینگ نیز حل‌شدنی است. برای ساخت ماشین‌های تورینگ
پیچیده‌تر، مهم است خود را قانع کنیم که می‌توانیم آن‌ها را ترکیب کنیم؛
در این صورت می‌توانیم با شکستن روش به بخش‌های ساده‌تر، ماشین‌هایی برای
حل مسائل پیچیده‌تر بسازیم. اگر راهی طبیعی برای تجزیهٔ مسئلهٔ پیچیده
به بخش‌های سازنده‌اش بیابیم، می‌توانیم در چند مرحله به مسئله بپردازیم:
چند ماشین تورینگ ساده بسازیم و آن‌ها را در ماشینی یگانه ترکیب کنیم
که مسئله را حل کند. این نکته هنگام پرداختن به مسئلهٔ توقف در بخش بعد
اهمیت ویژه‌ای دارد.

چگونه ماشین‌های تورینگ $M=\tuple{Q,\Sigma,q_0,\delta}$ و
$M'=\tuple{Q',\Sigma',q_0',\delta'}$ را ترکیب کنیم؟ پیکربندی نوار پس
از توقف~$M$ را به‌عنوان پیکربندی ورودیِ اجرای ماشین~$M'$ به کار
می‌بریم. برای به‌دست‌آوردن ماشین تورینگ یگانهٔ $M\frown M'$ که چنین
می‌کند، این گام‌ها را بردارید:
\begin{enumerate}
    \item همهٔ حالت‌های~$Q'$ از~$M'$ را دوباره شماره‌گذاری (یا
    برچسب‌گذاری) کنید تا $M$ و~$M'$ هیچ حالت مشترکی نداشته باشند
    ($Q\cap Q'=\emptyset$).
    \item حالت‌های $M\frown M'$ برابر $Q\cup Q'$ است.
    \item الفبای نوار برابر $\Sigma\cup\Sigma'$ است.
    \item حالت آغازین~$q_0$ است.
    \item تابع گذار، تابع $\delta''$ زیر است:
    \[\delta''(q,\sigma) =
    \begin{cases}
      \delta(q,\sigma) & \text{اگر $q \in Q$ و $\delta(q,\sigma)$ تعریف شده باشد}\\
      \delta'(q,\sigma) & \text{اگر $q \in Q'$}\\
      \tuple{q_0', \sigma, \TMstay} & \text{اگر $q \in Q$ و
      $\delta(q,\sigma)$ تعریف‌نشده باشد.}
    \end{cases}\]
\end{enumerate}
ماشین حاصل هنگامی که در حالت $q\in Q$ است از دستورهای~$M$ و هنگامی
که در حالت $q\in Q'$ است از دستورهای~$M'$ استفاده می‌کند. وقتی در
حالت $q\in Q$ نماد~$\sigma$ای را می‌خواند که $M$ برایش گذاری ندارد
(یعنی $M$ متوقف می‌شد)، وارد حالت آغازین~$M'$ می‌شود و محتویات نوار
و جای سر را همان‌گونه باقی می‌گذارد.

توجه کنید اگر ماشین~$M$ منضبط نباشد، نمی‌دانیم هنگام توقف~$M$ سر
نوار کجاست؛ پس در پیکربندی توقف~$M$ لازم نیست سر خانهٔ~$1$ را بخواند.
هنگام ترکیب ماشین‌ها باید این نکته را در نظر داشت.
\end{explain}

\begin{ex}
\emph{ترکیب ماشین‌ها:} ماشینی طراحی می‌کنیم که با ورودی متشکل از دو
بلوک $\TMstroke$ به طول $n$ و~$m$ آغاز می‌شود و با بلوک یگانه‌ای از
$2(m+n)$ عدد $\TMstroke$ روی نوار متوقف می‌شود. برای ساخت این ماشین
می‌توانیم دو ماشین آشنای جمع و دوبرابرکن را ترکیب کنیم. با ترسیم
نمودار حالت ماشین جمع آغاز می‌کنیم.
\[
\begin{tikzpicture}[->,>=stealth',shorten >=1pt,auto,node distance=2.8cm,
                    semithick]
  \tikzstyle{every state}=[fill=none,draw=black,text=black]

  \node[initial,state] (A)              {$q_0$};
  \node[state]         (B) [right of=A] {$q_1$};
  \node[state]         (C) [right of=B] {$q_2$};

  \path (A) edge node {\TMtrans{\TMblank}{\TMstroke}{\TMstay}} (B)
            edge [loop above] node {\TMtrans{\TMstroke}{\TMstroke}{\TMright}} (A)
        (B) edge [loop above] node {\TMtrans{\TMstroke}{\TMstroke}{\TMright}} (B)
            edge node {\TMtrans{\TMblank}{\TMblank}{\TMleft}} (C)
        (C) edge [loop above] node {\TMtrans{\TMstroke}{\TMblank}{\TMstay}} (C);
\end{tikzpicture}
\]
به‌جای توقف در حالت~$q_2$ می‌خواهیم کار را برای دوبرابرکردن خروجی
ادامه دهیم. به یاد آورید ماشین دوبرابرکن نخستین خط ورودی را پاک و دو
خط در خروجی جداگانه‌ای می‌نویسد. دستوری بیفزاییم تا مطمئن شویم سر
نوار نخستین خط خروجی ماشین جمع را می‌خواند.
\[
\begin{tikzpicture}[->,>=stealth',shorten >=1pt,auto,node distance=2.8cm,
                    semithick]
  \tikzstyle{every state}=[fill=none,draw=black,text=black]

  \node[initial,state] (A)              {$q_0$};
  \node[state]         (B) [right of=A] {$q_1$};
  \node[state]         (C) [right of=B] {$q_2$};
  \node[state]         (D) [below left of=C] {$q_3$};
  \node[state]         (E) [below left of=D] {$q_4$};

  \path (A) edge node {\TMtrans{\TMblank}{\TMstroke}{\TMstay}} (B)
            edge [loop above] node {\TMtrans{\TMstroke}{\TMstroke}{\TMright}} (A)
        (B) edge [loop above] node {\TMtrans{\TMstroke}{\TMstroke}{\TMright}} (B)
            edge node {\TMtrans{\TMblank}{\TMblank}{\TMleft}} (C)
        (C) edge node {\TMtrans{\TMstroke}{\TMblank}{\TMleft}} (D)
        (D) edge [loop left] node {\TMtrans{\TMstroke}{\TMstroke}{\TMleft}} (D)
            edge node[left, xshift=-2mm] {\TMtrans{\TMendtape}{\TMendtape}{\TMright}} (E);
\end{tikzpicture}
\]
اکنون دوبرابرکردن ورودی آسان است---تنها باید ماشین دوبرابرکن را به
حالت~$q_4$ وصل کنیم. برای این کار باید نام حالت‌های ماشین دوبرابرکن
را چنان تغییر دهیم که به‌جای~$q_0$ از~$q_4$ آغاز شوند؛ بدین‌ترتیب دو
حالت آغازین نخواهیم داشت. نمودار نهایی باید مانند
\olref{fig:combined} باشد.
\begin{figure}
\[
\begin{tikzpicture}[->,>=stealth',shorten >=1pt,auto,node distance=2.8cm,
                    semithick]
  \tikzstyle{every state}=[fill=none,draw=black,text=black]
  \node[initial,state] (A)              {$q_0$};
  \node[state]         (B) [right of=A] {$q_1$};
  \node[state]         (C) [right of=B] {$q_2$};
  \node[state]         (D) [below left of=C] {$q_3$};
  \node[state]         (E) [below left of=D] {$q_4$};
  \node[state]         (2) [right of=E] {$q_5$};
  \node[state]         (3) [right of=2] {$q_6$};
  \node[state]         (4) [below of=3] {$q_7$};
  \node[state]         (5) [left of=4]  {$q_8$};
  \node[state]         (6) [left of=5]  {$q_9$};

  \path (A) edge node {\TMtrans{\TMblank}{\TMstroke}{\TMstay}} (B)
            edge [loop above] node {\TMtrans{\TMstroke}{\TMstroke}{\TMright}} (A)
        (B) edge [loop above] node {\TMtrans{\TMstroke}{\TMstroke}{\TMright}} (B)
            edge node {\TMtrans{\TMblank}{\TMblank}{\TMleft}} (C)
        (C) edge node {\TMtrans{\TMstroke}{\TMblank}{\TMleft}} (D)
        (D) edge [loop left] node {\TMtrans{\TMstroke}{\TMstroke}{\TMleft}} (D)
            edge node[left, xshift=-2mm] {\TMtrans{\TMendtape}{\TMendtape}{\TMright}} (E)
    (E) edge node {\TMtrans{\TMstroke}{\TMblank}{\TMright}} (2)
    (2) edge [loop above] node {\TMtrans{\TMstroke}{\TMstroke}{\TMright}} (2)
      edge node {\TMtrans{\TMblank}{\TMblank}{\TMright}} (3)
    (3) edge [loop above] node {\TMtrans{\TMstroke}{\TMstroke}{\TMright}} (3)
        edge node {\TMtrans{\TMblank}{\TMstroke}{\TMright}} (4)
    (4) edge [loop below] node {\TMtrans{\TMblank}{\TMstroke}{\TMleft}} (4)
        edge node {\TMtrans{\TMstroke}{\TMstroke}{\TMleft}} (5)
    (5) edge [loop below]  node {\TMtrans{\TMstroke}{\TMstroke}{\TMleft}} (5)
        edge              node {\TMtrans{\TMblank}{\TMblank}{\TMleft}} (6)
    (6) edge [loop below] node {\TMtrans{\TMstroke}{\TMstroke}{\TMleft}} (6)
        edge              node {\TMtrans{\TMblank}{\TMblank}{\TMright}} (E);
\end{tikzpicture}
\]
\caption{ترکیب ماشین‌های جمع و دوبرابرکن}
\ollabel{fig:combined}
\end{figure}
\end{ex}

\begin{prop}
    اگر $M$ و~$M'$ منضبط باشند و به‌ترتیب توابع
    $f\colon\Nat^k\to\Nat$ و $f'\colon\Nat\to\Nat$ را محاسبه کنند،
    آنگاه $M\frown M'$ منضبط است و~$\comp{f}{f'}$ را محاسبه می‌کند.
\end{prop}

\begin{proof}
    چون $M$ منضبط است، هنگامی که با خروجی
    $f(n_1,\dots,n_k)=m$ متوقف می‌شود، سر خانهٔ~$1$ را می‌خواند. اگر
    اکنون وارد حالت آغازین~$M'$ شویم، ماشین با خروجی $f'(m)$ و باز هم
    درحالی‌که خانهٔ~$1$ را می‌خواند متوقف می‌شود. دیگر شرط‌های
    \olref[dis]{defn:disciplined} نیز برقرارند.
\end{proof}

\begin{prob}
    با گرفتن ماشین~$M$ از \cref{tur:mac:dis:prob:disc-succ} و ساختن
    $M\frown M$، ماشین تورینگ منضبطی به دست دهید که $f(x)=x+2$ را
    محاسبه کند.
\end{prob}
\end{document}
